% !PRINTED 759
\weberpage{759}
%<<CONT
is an integral rational number divisible by $p$. From this it follows finally that
\[
CF(0) + F(\alpha_1) + F(\alpha_2) + \cdots + F(\alpha_\mu)
\]
is an integral rational number not divisible by $p$, and therefore also not vanishing, which consequently, apart from its sign, is at least $=1$.

If we can now finally prove that, under the present assumption about $f$ as well, the function $Q(x)$ can be made arbitrarily small for every finite $x$ by making $p$ sufficiently large, then the impossibility of equation (4) follows exactly as above.

But this may be seen in the following way: instead of the function
\weq{6}{f(x) = c_nx^n + c_{n-1}x^{n-1} + \cdots + c_0}
we consider the function
\[
\begin{aligned}
\psi(x) &= \dfrac{a^{\mu p+p-1}x^{p-1}(x+a_1)^p(x+a_2)^p\ldots(x+a_\mu)^p}{\Pi(p-1)} \\
&= \gamma_nx^n + \gamma_{n-1}x^{n-1} + \cdots + \gamma_0,
\end{aligned}
\]
which now has none but positive (though not necessarily rational) coefficients. The coefficients $c_n,c_{n-1}\ldots$ are formed by multiplication and addition from the numbers $a,-\alpha_1,-\alpha_2,\ldots,-\alpha_\mu$, and the corresponding coefficients $\gamma_n,\gamma_{n-1}\ldots$ are obtained from them if the $-\alpha_1,-\alpha_2,\ldots,-\alpha_\mu$ are replaced by their absolute values $a_1,a_2,\ldots,a_\mu$; from which it follows, by the theorem of the Introduction already mentioned above, that the coefficients $\gamma_n,\gamma_{n-1},\ldots$ are certainly not smaller than the absolute values of the corresponding $c_n,c_{n-1}\ldots$

Now for every finite $x$ whose absolute value is $\xi$, the absolute value of $Q(x)$ is, by \S.\,205, (10), smaller than, or at least not greater than,
\[
\gamma_n\xi^n + \gamma_{n-1}\xi^{n-1} + \cdots + \gamma_0 = \psi(\xi),
\]
and that $\psi(\xi)$ sinks below every limit when $p$ becomes large enough is inferred from the representation:
\[
\psi(x) = \dfrac{X}{ax}\dfrac{X^{p-1}}{\Pi(p-1)},
\]
if
\[
% [sic] -- the source prints Latin $u+1$; the derivation from the exponent $\mu p+p-1$ requires $\mu+1$.
X = a^{u+1}x(x+a_1)(x+a_2)\ldots(x+a_\mu)
\]
is put.
